\documentclass[11pt, a4paper]{article}

% --- REQUIRED SYSTEM PREAMBLE ---
\usepackage[a4paper, margin=2.2cm, top=2.5cm, bottom=2.5cm]{geometry}
\usepackage{fontspec}
\usepackage[english, vietnamese, bidi=basic, provide=*]{babel}
\babelprovide[import, main, onchar=ids fonts]{vietnamese}
\babelfont{rm}{Noto Sans}

% --- MATH & FORMATTING ---
\usepackage{amsmath, amssymb, mathtools}
\usepackage{microtype}
\usepackage{parskip}

% --- COLORS & BOXES ---
\usepackage{xcolor}
\definecolor{biomedblue}{RGB}{27,75,138}
\definecolor{biomeddark}{RGB}{15,32,67}
\definecolor{accentteal}{RGB}{0,128,128}
\definecolor{boxgray}{RGB}{245,247,250}
\definecolor{bordergray}{RGB}{210,215,225}
\definecolor{alertred}{RGB}{217,83,79}

\usepackage[most]{tcolorbox}
\newtcolorbox{infobox}[1][]{
  enhanced, boxrule=1pt, arc=3pt, left=10pt, right=10pt, top=8pt, bottom=8pt,
  colback=biomedblue!4!white, colframe=biomedblue!85!black,
  title={#1}, coltitle=white, fonttitle=\bfseries
}

% --- TABLES & CODE BLOCKS ---
\usepackage{booktabs}
\usepackage{tabularx}
\usepackage{listings}
\lstset{
  basicstyle=\ttfamily\footnotesize,
  backgroundcolor=\color{boxgray},
  frame=single,
  rulecolor=\color{bordergray},
  breaklines=true,
  showstringspaces=false,
  tabsize=2
}

% --- HEADERS & FOOTERS ---
\usepackage{fancyhdr}
\pagestyle{fancy}
\fancyhf{}
\fancyhead[L]{\small\color{gray}\textsc{HUST SOICT -- IT5425: Quản trị Dữ liệu \& Trực quan hóa}}
\fancyhead[R]{\small\color{gray}\textsc{Báo cáo Phân hệ 3 -- TV 3 (EDA)}}
\fancyfoot[C]{\thepage}
\renewcommand{\headrulewidth}{0.4pt}
\renewcommand{\footrulewidth}{0.4pt}

% Hyperref must be last
\usepackage[colorlinks=true,linkcolor=biomedblue,citecolor=biomedblue,urlcolor=biomedblue]{hyperref}

\title{%
  \vspace{-1.5cm}
  {\Large\bfseries TRƯỜNG CÔNG NGHỆ THÔNG TIN VÀ TRUYỀN THÔNG -- ĐẠI HỌC BÁCH KHOA HÀ NỘI}\\[0.3cm]
  {\large HỌC PHẦN: QUẢN TRỊ DỮ LIỆU VÀ TRỰC QUAN HÓA (IT5425)}\\[0.6cm]
  \rule{\textwidth}{1.5pt}\\[0.3cm]
  {\LARGE\bfseries BÁO CÁO KỸ THUẬT TIẾN ĐỘ PHÂN HỆ 3}\\[0.2cm]
  {\large\color{biomedblue}\textbf{Phân Tích Khám Phá Dữ Liệu Y Sinh (EDA), Kiểm Định Giả Thuyết \& Trực Quan Hóa Tỷ Số Chênh (Odds Ratios)}}\\[0.2cm]
  \rule{\textwidth}{0.8pt}
}
\author{%
  \textbf{Thành viên 3: EDA \& Biostatistics}\\[0.1cm]
  Nhóm Nghiên cứu Y sinh (Biomedicine Research Group) -- Lớp IT5425
}
\date{Tháng 9 Năm 2026}

\begin{document}
\maketitle
\thispagestyle{fancy}

\begin{abstract}
Tài liệu này tổng kết toàn bộ kết quả triển khai kỹ thuật của \textbf{Phân hệ 3 (TV 3 -- EDA)} trong dự án Bài tập lớn học phần \textit{Quản trị Dữ liệu và Trực quan hóa (IT5425)}. Phân hệ 3 chịu trách nhiệm tiếp nhận dữ liệu đã được làm sạch từ kho LanceDB, tiến hành quy trình Phân tích Khám phá (EDA) và kiểm định giả thuyết thống kê chuyên sâu trên tập mẫu $\mathbf{8{,}828}$ hồ sơ bệnh nhân đủ điều kiện. Hệ thống tự động hóa toàn bộ việc xây dựng Bảng đặc điểm lâm sàng (Table 1), áp dụng thuật toán quyết định kiểm định phi tham số (\textit{Mann-Whitney U}, \textit{Chi-square}), và đánh giá mức độ rủi ro mạch vành 10 năm qua mô hình hồi quy Logistic (\textit{Odds Ratios}). Điểm nhấn kỹ thuật là việc chuẩn hóa Z-score các biến liên tục, trực quan hóa Split-Violin và Forest Plot chuẩn học thuật, cùng chiến lược sàng lọc đặc trưng tinh gọn để chuyển giao cho các phân hệ phân tích dự báo.
\end{abstract}

\vspace{0.3cm}
\tableofcontents
\vspace{0.5cm}

% ==============================================================================
\section{Tổng Quan Mục Tiêu \& Vị Trí Cốt Lõi Của Phân Hệ 3}
% ==============================================================================

Trong hệ sinh thái luồng dữ liệu của dự án, Phân hệ 3 đóng vai trò "người giải mã" các tín hiệu lâm sàng trước khi đưa vào huấn luyện mô hình học máy:
\begin{itemize}
    \item \textbf{Cung cấp Insight Lâm Sàng:} Đánh giá sự khác biệt phân phối của các chỉ số sinh hóa (tuổi, huyết áp, đường huyết, BMI, cholesterol) giữa nhóm có và không có biến cố tim mạch.
    \item \textbf{Cung cấp Báo cáo Học thuật (Table 1):} Tự động tổng hợp và tính toán báo cáo đặc điểm bệnh nhân ban đầu, tự động chuyển đổi sang các đại lượng thống kê phù hợp như Median [IQR] cho dữ liệu không chuẩn.
    \item \textbf{Sàng Lọc Đặc Trưng (Feature Selection) cho Phân hệ Học máy:} Bảng phân tích Tỷ số chênh (Odds Ratios) là căn cứ định lượng Toán học vững chắc để định hướng cho việc xây dựng mô hình rủi ro tại Phân hệ 4 và 5.
\end{itemize}

% ==============================================================================
\section{Đặc Tả Kỹ Thuật Các Tệp Tin Mã Nguồn Đã Xây Dựng}
% ==============================================================================

Phân hệ 3 đã xây dựng và chuẩn hóa 3 cấu phần mã nguồn phục vụ tính toán thống kê:

\subsection{Thư viện Kiểm định Thống kê \& Biểu đồ (\texttt{src/biomed\_lake/biostats/test.py})}
Module chịu trách nhiệm phân tích thống kê mô tả và suy diễn (Univariate Analysis):
\begin{itemize}
    \item \textbf{\texttt{generate\_table\_one}:} Tích hợp thư viện \texttt{tableone}, hỗ trợ tham số \texttt{nonnormal\_cols} để xuất báo cáo dưới dạng Trung vị [Khoảng tứ phân vị] thay vì Mean (SD) đối với dữ liệu bị lệch phân phối. Xuất tự động định dạng CSV và HTML.
    \item \textbf{\texttt{run\_hypothesis\_test}:} Thuật toán kiểm định động thông minh tự động rẽ nhánh sử dụng \textbf{Mann-Whitney U test} cho biến liên tục vi phạm giả định chuẩn, và \textbf{Chi-square test} cho biến phân loại.
    \item \textbf{\texttt{plot\_biochemical\_distributions}:} Vẽ trực quan hóa bằng kỹ thuật Split-Violin plot lồng ghép Boxplot để biểu diễn mật độ xác suất đối sánh trực tiếp 2 nhóm bệnh lý.
\end{itemize}

\subsection{Thư viện Mô hình Logistic \& Tỷ Số Chênh (\texttt{src/biomed\_lake/biostats/odds.py})}
Module đánh giá rủi ro đa biến (Multivariate Analysis):
\begin{itemize}
    \item \textbf{\texttt{calculate\_or\_ci}:} Áp dụng chuẩn hóa \textbf{StandardScaler (Z-score)} cho các biến liên tục. Điều này giúp Tỷ số chênh (OR) phản ánh mức tăng rủi ro tương ứng với $1$ độ lệch chuẩn ($\sigma$), mang lại góc nhìn tương quan công bằng với các biến phân loại. Mô hình được xây dựng qua \texttt{sm.Logit} để trích xuất khoảng tin cậy 95\% (CI).
    \item \textbf{\texttt{plot\_forest}:} Thuật toán vẽ Forest Plot trực quan hóa các khoảng tin cậy 95\% và Tỷ số chênh (OR). Biểu đồ được thiết kế chuẩn mực học thuật, tập trung hiển thị độ lớn của rủi ro, tuân thủ nguyên tắc tính nguyên vẹn của kết quả thống kê.
\end{itemize}

\subsection{Luồng Thực Thi Toàn Trình (\texttt{pipelines/03\_eda.py})}
Kịch bản Pipeline kết nối kho lưu trữ LanceDB, trích xuất Polars DataFrame, kết nối các bảng chiều (Dimension) và bảng thực tế (Fact), loại bỏ dữ liệu khuyết thiếu và kích hoạt song song chuỗi phân tích. Điểm nhấn là việc chuyển đổi thành công cột \texttt{bp\_stage} thành dữ liệu Categorical có thứ tự, tối ưu hóa quá trình nội suy quy chiếu của thuật toán.

% ==============================================================================
\section{Kết Quả Phân Tích Thực Nghiệm}
% ==============================================================================

\subsection{Kết Quả Chạy Pipeline 03}
Hệ thống xử lý $\mathbf{8{,}828}$ bệnh nhân đủ điều kiện và in ra báo cáo thống kê:

\begin{lstlisting}[language=bash]
================================================================================
[*] PIPELINE 03: BIOMEDICAL EDA & HYPOTHESIS TESTING
================================================================================
[*] Dang lay du lieu tu LanceDB
[*] Kich thuoc mau du dieu kien phan tich: 8828 benh nhan

--- Giai doan 1: Xay dung Table 1 ---
[+] Dang xay dung Table 1

--- Giai doan 2: Kiem dinh gia thuyet va ve phan phoi ---
[+] Dang kiem dinh gia thuyet
[Ket qua kiem dinh gia thuyet]
          Bien           Kiem dinh       P-value
           age Mann-Whitney U test 2.618994e-159
        sys_bp Mann-Whitney U test 3.320882e-154
        dia_bp Mann-Whitney U test 1.304637e-106
      tot_chol Mann-Whitney U test  3.512924e-05
           BMI Mann-Whitney U test  2.655625e-83
       glucose Mann-Whitney U test 1.109548e-148
       is_male          Chi-square  8.196966e-01
current_smoker          Chi-square  6.420470e-50
      bp_stage          Chi-square 1.814963e-206
      diabetes          Chi-square  1.742776e-27
       bp_meds          Chi-square  6.457047e-01
[+] Dang ve bieu do Split-Violin va Box cho cac chi so sinh hoa

--- Giai doan 3: Tinh toan ti so chenh lech OR va ve bieu do ---
[+] Dang tinh toan Odds Ratio
[Bang OR]
                               CI_lower  CI_upper        OR       p_value
age                            1.554437  1.734007  1.641769  1.065646e-70
tot_chol                       0.780185  0.870939  0.824314  5.886833e-12
BMI                            1.199175  1.327871  1.261883  3.748574e-19
glucose                        1.549288  1.859195  1.697183  5.815642e-30
is_male                        1.092255  1.348847  1.213789  3.190727e-04
current_smoker                 0.617207  0.804815  0.704796  2.378338e-07
diabetes                       0.283383  0.569040  0.401567  2.895679e-07
bp_meds                        0.317220  0.775254  0.495909  2.093361e-03
bp_stage_elevated              0.980462  1.660400  1.275915  6.981020e-02
bp_stage_stage_1_hypertension  1.362991  1.898721  1.608708  1.888532e-08
bp_stage_stage_2_hypertension  4.032167  5.653561  4.774526  1.808858e-73
bp_stage_hypertensive_crisis   3.514825  6.746473  4.869566  1.785337e-21
[+] Dang ve bieu do Forest
[OK] PIPELINE 3 HOAN TAT
\end{lstlisting}

\subsection{Phân Tích Đơn Biến \& Đa Biến}
\begin{itemize}
    \item \textbf{Phân tích đơn biến:} Hầu hết các biến số sinh hóa và lâm sàng (tuổi, huyết áp tâm thu/tâm trương, BMI, glucose) đều cho thấy sự phân tách rủi ro cực kỳ mạnh mẽ giữa nhóm bệnh và nhóm khỏe mạnh ($p$-value $\sim 0$).
    \item \textbf{Hồi quy Logistic Đa biến:} Kết quả chỉ ra bệnh nhân thuộc nhóm \texttt{hypertensive\_crisis} và \texttt{stage\_2\_hypertension} có rủi ro mắc bệnh mạch vành (OR) tăng vọt lên tới $\mathbf{4.86}$ và $\mathbf{4.77}$ lần. Kỹ thuật chuẩn hóa cũng làm nổi bật tác động của sự gia tăng đường huyết và tuổi tác (tăng rủi ro thêm khoảng 64\%-69\% trên mỗi độ lệch chuẩn).
\end{itemize}

\end{document}